% 附录：冻结符号预测与可计算证书（完整材料；修订版用）。
% 位置：现有附录之后。不进正文、摘要、贡献列表。
% 口径与 signcriterion_note.md §13–§15 一致（终版口径）。

\section{Frozen Sign Prediction and Computable Certificates}
\label{app:signpred}

This appendix reports, as further validation of the main line, a
test of how far the mechanistic decomposition supports prediction.
All sign rates use the 88-point sweep grid ($2$ noise rates
$\times$ $4$ depths $\times$ $11$ coupling values), restricted to
$|\Delta|>10^{-3}$; $\Delta$ denotes the mean-margin difference
$\bar m^{A}_T-\bar m^{B}_T$ under the matched-marginal controls
of Sec.~2.

\paragraph{Protocol.}
The protocol has three frozen stages: calibration (rules fitted on
$p\le0.15$, $T\le10$, coupling endpoints $\varphi\in\{0,1\}$),
prediction (written to disk before any evaluation run; disjoint
seed ranges), and evaluation. The coupling $\varphi$ interpolates
between truth-initialized ($\varphi{=}0$) and corruption-
initialized ($\varphi{=}1$) readouts: with probability $\varphi$,
a corrupted node carries the same corrupted label that generated
the $A$-graph. Four settings are studied: Cora and DBLP under the
heat filter, and each under the normalized APPNP filter
($S=D^{-1/2}(A{+}I)D^{-1/2}$, $Z\leftarrow0.9\,SZ+0.1\,Z_0$);
DBLP/APPNP is excluded from all rule and constant selection and
serves as the confirmation setting.

\paragraph{Rules.}
Both rule variants use the exact decomposition
$E[\Delta]=p\,G+(1-p)\,D$ with the isolation term $1-2\varphi$
(closed form) and a fitted recovery curve; they differ only in
whether the clean contrast is mixed in $\varphi$ (v2: $D(\varphi)$
interpolated between calibrated endpoints; v1: $D\equiv D_{0}$).
Table~\ref{tab:app-rules} reports margin-sign accuracy together
with constant-sign and development-majority baselines.

\begin{table}[h]
\centering\small
\begin{tabular}{lcccccc}
\hline
setting & pos/neg & const-neg & dev-majority & v1 & v2 \\
\hline
Cora/heat      & 12/60 & 83\% & 83\% & 99\%  & 100\% \\
DBLP/heat      & 20/63 & 76\% & 76\% & 86\%  & 98\%  \\
Cora/APPNP     & 12/56 & 82\% & 82\% & 94\%  & 96\%  \\
DBLP/APPNP (conf.) & 0/63 & 100\% & 100\% & 100\% & 49\% \\
\hline
\end{tabular}
\caption{Sweep-grid sign accuracy with baselines. The
confirmation setting contains no positive configurations: a
constant-negative predictor attains 100\%, so it tests only
negative-direction prediction; the unmixed variant's advantage
over constant-negative is demonstrated on the development
settings (positive-direction accuracy 95--100\%) and its
extrapolation is not established. Whether the mixed variant's
failure reflects in-family overfitting or component mismatch in
this regime is not distinguished.}
\label{tab:app-rules}
\end{table}

\paragraph{Component-level facts.}
Holding uniformly across operators and datasets: (i) conditional
on isolation under $A$,
$E[m^A_T(U)\mid W_U{=}1,\ U\ \text{isolated}]=1-2\varphi$
exactly---the protocol draws the initialization-coupling coin
independently of the corruption realization, so
$\Pr[\text{init}\ne y\mid W_U{=}1,\text{isolated}]=\varphi$ and
an isolated readout is preserved; the unconditional expectation
carries a remainder bounded by an isolation-probability-weighted
displacement term (measured deviation $\le0.026$, an empirical
observation); (ii) the recovery curve $E[m^B_T\mid W_U{=}1]$ is
linear in $\varphi$ on all four settings.

\paragraph{Two-sided certificates.}
Replacing fitted terms by instance-computable ones (closed-form
retention rates, oracle-graph response curves, expectation-
operator mean field, A-only/B-only simulation audits) yields
sufficient conditions in both directions:
$p\,G^{lb}+(1-p)(\hat D-\nu)>0\Rightarrow E[\Delta]>0$, and the
mirror inequality for $E[\Delta]<0$. Three layers are distinct:
the implication is mathematical; the model-error slack $\nu$ is
an assumption audited on single-model realizations (in the
confirmation stage no component reads the A--B difference, while
the slack constants were fixed from development-stage results
that did use difference-based diagnostics); Monte-Carlo
components carry asymptotic one-sided bounds ($z=3.09$). The grid
involves $42$ one-sided tail events ($\beta_c,\beta_w$
two-sided across $16$ cells; $\bar\varepsilon$ and $E[\delta^2]$
upper-only); their Bonferroni asymptotic nominal failure level is
about $4.2\%$, and this concerns the MC layer only---certificate
validity separately requires the mathematical premises and, at
one tier, the model-error assumption. No finite-sample claim is
made. On evaluated instances no false certificate was observed
(abstention included): coverage is $88\%$ on Cora/heat at the
$(S')$-conditional tier and $38\%$ at the $(S')$-free asymptotic
conditional tier, $13\%$ on DBLP/heat, $15\%$ on Cora/APPNP, and
the frozen criterion issues no certificates on DBLP/APPNP---an
abstention that exposes the coverage limitation rather than
evidencing soundness.

\paragraph{Budget comparison.}
Under a reuse-cost model (one operator pair and eleven flows per
realization cover the sweep grid; depths recorded en route),
budget-matched direct Monte Carlo shows no demonstrated
certificate advantage on Cora/heat ($82\%$ vs $88\%$; bootstrap
CI on the difference $[-0.02,+0.13]$) and higher coverage on the
other three settings ($100\%$ vs $13\%$; $77\%$ vs $15\%$;
$100\%$ vs $0\%$).

\paragraph{Correction log and reproducibility.}
The development history includes three audit-driven corrections
(a certificate inequality using the wrong one-sided bound; a
mixed-edge retention formula; a sign error in the recovery
correction), each exposed by independent recomputation and
documented in the lab note; all reported numbers are reproducible
from the frozen scripts (\texttt{sign\_predict.py},
\texttt{sign\_protocol.py}, \texttt{mc\_baseline.py},
\texttt{baseline\_fair.py}) with independent audit scripts
verifying certificate recomputation, scoring, seed disjointness,
and cross-validation against the paper's original implementation
(12 endpoint comparisons in agreement).
